\subsection{Design Space}

Based on the analysis, we derived a hierarchical design space that organizes strategies for designing failure across the three facets: Failure Emergence, Failure Experience, and Failure Response. Each facet contains multiple dimensions, for a total of seven, representing specific characteristics within that facet. At the lowest level, concrete design decisions operationalize each dimension. Table.~\ref{tab:design space} summarizes the design space, with check marks indicating the presence of a particular design decision in a given project.

\subsubsection{Failure Emergence}

Failure emergence refers to how failure arises inherently from data structures and system rules. In this process, designers need to consider what kinds of \textit{structural constraints} characterize failure, and under what \textit{conditions} it is \textit{triggered}.

\paragraph{Structural Constraints.}
This dimension describes how the data, models, and rules constrain possible outcomes, making certain goals difficult or impossible to achieve.

Data may lead to \textbf{pre-determined outcomes} (n=11), which refer to cases where the outcome of failure is already determined by the underlying data or model before the game begins. Player actions cannot alter the final outcome, but can only gradually approach or reveal this predetermined result. For instance, \textit{Inevitable Super Rich}~\cite{InevitableSuperRich} (Fig.~\ref{fig:example}-A) allows players to conduct random coin-flip trades with the rich and observe changes in wealth distribution, demonstrating the process of wealth concentration in a free market. While players may attempt to gain wealth from the rich, the underlying model amplifies wealth inequality, rendering such outcomes unattainable and leading to an inevitable concentration of wealth among a small number of individuals, resulting in failure.

Data may impose \textbf{bounded value ranges} (n=13), where numerical constraints such as limited total resources or capped indicators restrict players to optimizing within a limited range and prevent them from achieving success. For example, in \textit{The En-ROADS}~\cite{TheEn-ROADS} (Fig.~\ref{fig:example}-B), players simulate adjusting parameters such as energy structure, carbon emissions, and technology investment proportions to mitigate global warming. Due to real-world resource constraints and efficiency limits, players cannot adjust these parameters to their theoretical optimum to halt global warming.

Data may introduce \textbf{conflicting metrics} (n=14), where variables are interdependent and impose trade-offs, such that improvements in one metric come at the expense of others and multiple metrics cannot simultaneously reach their ideal states. For example, in \textit{Educate Your Child}~\cite{EducateYourChild} (Fig.~\ref{fig:educate}), players attempt to select a school that performs well across all criteria. However, negative correlations between academic performance and racial diversity, as well as between household income and proximity to school, make it difficult to satisfy all of these criteria at the same time.

\paragraph{Triggering Conditions.}
This dimension describes the circumstances under which failure occurs, resulting from interactions between system dynamics and player actions.

Failure may be \textbf{decision-driven} (n=20), referring to cases where suboptimal player decisions lead to failure, such as ineffective strategies or inappropriate sequencing of actions. For example, in \textit{The Nest Egg Game}~\cite{TheNestEggGame} (Fig.~\ref{fig:example}-C), players make financial decisions across different life stages. Poor savings or investment choices may result in debt and eventual failure, whereas adjusting decision strategies can improve outcomes.

Failure may also be \textbf{knowledge-constrained} (n=11), referring to cases where players lack sufficient background knowledge or an adequate understanding of system rules when making decisions. For example, in \textit{Salubrious Guess}~\cite{SalubriousGuess} (Fig.~\ref{fig:example}-D), players estimate the values of health indicators in a randomly selected region based on their existing knowledge. As players acquire relevant knowledge, they can make more accurate estimates and achieve higher scores.

\begin{figure*}[!t]
  \centering
  \includegraphics[width=\textwidth]{figures/example.png}
  \caption{Example {\name} projects from our corpus. Example sources: ~\cite{InevitableSuperRich,TheEn-ROADS,TheNestEggGame,SalubriousGuess,USPresidentialElectiony,CollegeChancesbyIncome,AbortionMazes,BirthControlMeasures,ToBuildaBetterBallot,HeartSaver}}
  \label{fig:example}
\end{figure*}

Failure can be \textbf{externally-driven} (n=1), referring to cases where failure arises from factors beyond the player's control, such as the actions of others, particularly in collaborative contexts. For example, in \textit{People of the Pandemic}~\cite{PeopleofthePandemic} (Fig.~\ref{fig:pandemic}), players play alongside a team of 20 real people. Even if a player chooses to stay home almost entirely for a week, the spread of infection cannot be quickly reduced because the player cannot control how often others go out.

Finally, failure can be \textbf{systemic} (n=6), referring to cases where certain goals remain difficult to achieve even with well-reasoned decisions due to inherent structural constraints of the system. For example, in \textit{Inevitable Super Rich}~\cite{InevitableSuperRich} (Fig.~\ref{fig:example}-A), players find that regardless of their strategies, wealth distribution gradually concentrates among a small number of individuals, as the underlying Yard-sale model inherently produces this outcome.

\subsubsection{Failure Experience}

Failure experience refers to how systems represent failure in perceptible and interpretable forms, enabling players to perceive constraints and understand failure outcomes. In this process, designers need to consider how to \textit{present} pre-failure \textit{constraints} and how to \textit{encode failure}.

\paragraph{Constraints Presentation.}
This dimension describes when and how system constraints are exposed through visualization, allowing players to perceive the conditions under which failure occurs.

Constraints may be presented through \textbf{explicit visual marks} (n=4), which describe cases where constraints are directly represented through visualization elements such as threshold lines, upper bounds, or warning indicators, allowing players to clearly perceive them prior to interaction. For example, in \textit{US Presidential Election}~\cite{USPresidentialElectiony} (Fig.~\ref{fig:example}-E), the threshold of electoral votes required for victory is explicitly marked, allowing players to adjust the candidate's vote share accordingly to surpass this threshold.

Constraints may be presented through \textbf{implicit indicators} (n=14), which capture cases where failure conditions are not directly represented, but are indirectly reflected through changes in data or indicators, such as curve trends, slowing growth, or diminishing effects, allowing players to gradually perceive these constraints during interaction. For example, in \textit{The En-ROADS}~\cite{TheEn-ROADS} (Fig.~\ref{fig:example}-B), player decisions are transformed into temperature prediction curves, through which players interpret the effects and limitations of their adjustments.

Constraints may also be presented through \textbf{post-hoc consequences} (n=17), which occur when the system does not explicitly present constraints in advance, but instead reveals them when certain conditions are met through outcomes such as game termination or revealing the correct answers. For example, in \textit{College Chances by Income}~\cite{CollegeChancesbyIncome} (Fig.~\ref{fig:example}-F), players first draw a perceived relationship between income and college attendance probability. Upon completion, the system reveals real data, allowing players to recognize discrepancies between their subjective perceptions and actual data.

\paragraph{Failure Encoding.}
This dimension describes how failure is encoded into perceivable and interpretable outcomes that can be recognized by players.

Failure can be encoded as \textbf{comparative performance} (n=15), which describes cases where failure is not treated as an independent outcome, but is instead reflected through the relative gap between the player and a target or other players, such as scores, rankings, or progress, allowing players to perceive their distance from success within a comparative frame. For example, in \textit{Salubrious Guess}~\cite{SalubriousGuess} (Fig.~\ref{fig:example}-D), scores are assigned based on the deviation between guessed values and actual health indicators, and players can compare their performance on a leaderboard after gameplay.

Failure can be encoded as \textbf{evaluative feedback} (n=21), which describes cases where failure functions as immediate feedback on the results of current actions, signaling unmet goals and guiding players to adjust their strategies based on these outcomes, thereby gradually improving their behavior. For example, in \textit{51 Sprints}~\cite{51sprints} (Fig.~\ref{fig:sprints}), after adjusting factors and running the simulation, players determine failure based on whether the selected athlete wins or not, and continuously refine their strategies according to this feedback until the athlete wins.

Failure can also be encoded as \textbf{difficulty mapping} (n=3), which captures cases where the scale and frequency of failure are treated as quantifiable degrees, often reflected through the time and effort players must expend to overcome repeated setbacks. In this way, failure is used to represent the difficulty or complexity of the underlying problem, allowing players to perceive the severity of the issue. For example, in \textit{Abortion Mazes}~\cite{AbortionMazes} (Fig.~\ref{fig:example}-G), the author quantifies real policy data on abortion restrictions across states into complexity scores and generates mazes of corresponding difficulty, such that more difficult mazes correspond to more restrictive abortion policies, allowing players to perceive the severity of these policies.

\subsubsection{Failure Response}

Failure response refers to how players react to failure and take subsequent actions to adjust, understand, and reassess it. In this process, designers need to consider the multiple levels of response that players may adopt, including \textit{action adjustment}, \textit{exploratory learning}, and \textit{interpretive reframing}.

\paragraph{Action Adjustment.}
This dimension describes how players use their agency to make adjustments based on their existing understanding and attribution, in order to attempt to improve outcomes.

Players can perform \textbf{parameter tuning} (n=12), which involves improving outcomes by adjusting numerical variables or system parameters, such as changing resource proportions, policy intensity, or other controllable parameters, thereby optimizing current performance. For example, in \textit{Birth Control Measures}~\cite{BirthControlMeasures} (Fig.~\ref{fig:example}-H), players adjust fertility rates and policy intensity across regions and observe their impact on the global population size in 2050, thereby attempting to approach the goal of reducing population.

Players can also change \textbf{scope \& filtering selection} (n=2), which involves adjusting decisions by changing the subset of data under consideration, such as selecting different regions, groups, or time intervals, or applying different filtering conditions to obtain diverse results. For example, in \textit{To Build a Better Ballot}~\cite{ToBuildaBetterBallot} (Fig.~\ref{fig:example}-I), players switch voting rules or voter groups and compare multiple scenarios to explore improved voting outcomes.

Players can also engage in \textbf{action sequence modification} (n=5), which involves exploring alternative strategies by revising the timing or order of actions, such as rearranging policy implementation schedules or adjusting action timing to pursue different effects. For example, in \textit{Heart Saver}~\cite{HeartSaver} (Fig.~\ref{fig:example}-J), players adjust the order and distance priority of transporting patients, balancing distance and hospital quality while accounting for declining survival rates over time.

Players can also engage in \textbf{strategy reconfiguration} (n=5), which involves exploring possible solutions by combining different strategies or options, and optimizing decisions by comparing outcome differences across various combinations. For example, in \textit{The Nest Egg Game}~\cite{TheNestEggGame} (Fig.~\ref{fig:example}-C), players combine financial decisions across life stages to explore which arrangements are more effective in achieving their financial goals.

\paragraph{Exploratory Learning.}
This dimension describes how players gradually develop a deeper understanding of failure through interaction with the system, enabling them to identify valid outcomes as well as the underlying rules and patterns of the system.

Players can engage in \textbf{iterative experimentation} (n=8), which involves exploring different pathways through repeated attempts and continuous experimentation, thereby gradually accumulating an understanding of system behavior, variable relationships, and the issue itself. For example, in \textit{Heart Saver}~\cite{HeartSaver} (Fig.~\ref{fig:example}-J), players repeatedly adjust transportation routes and hospital selection, gradually understanding the relationships among time, distance, and hospital quality, as well as the dynamics of how survival rates change over time.

Players can also engage in \textbf{hypothesis testing} (n=10), which involves approaching reasonable decisions by estimating or inferring possible outcomes under conditions of incomplete information, while helping players develop familiarity with relevant domain knowledge. For example, in \textit{College Chances by Income}~\cite{CollegeChancesbyIncome} (Fig.~\ref{fig:example}-F), players first draw a perceived relationship between income and college admission probability, after which the system displays real data, allowing players to revise their understanding through comparison.

\paragraph{Interpretive Reframing.}
This dimension describes the process by which players define and reinterpret the causes and meanings of failure after forming an understanding of it.

Players can engage in \textbf{perspective shifting} (n=8), which involves reinterpreting the problem by shifting the represented role or standpoint, reattributing failure to different groups or situational conditions, and reassessing failure from a new perspective. For example, in \textit{Educate Your Child}~\cite{EducateYourChild} (Fig.~\ref{fig:educate}), players initially choose a child from a family with specific racial and income characteristics and act as their parent. From this perspective, they make decisions for the child, thereby experiencing the structural constraints faced by different groups.

Players can also engage in \textbf{criteria reweighting} (n=5), which involves interpreting failure as arising from conflicts among multiple goals, and thus adjusting the priority among different objectives to redefine the criteria for what counts as success, thereby reshaping the meaning of failure. For example, in \textit{51 Sprints}~\cite{51sprints} (Fig.~\ref{fig:sprints}), some players no longer pursue speed or medals, but instead treat building a fair environment unaffected by gender, class, and race as a new goal.
